\documentclass[11pt]{article}
\usepackage[a4paper, margin=2.5cm]{geometry}
\usepackage[utf8]{inputenc}
\usepackage[T1]{fontenc}
\usepackage[slovene]{babel}
\usepackage{amsmath,amssymb,amsthm}
\usepackage{graphicx}
\usepackage{hyperref}
\newtheorem{definition}{Definicija}
\newtheorem{theorem}{Izrek}
% Brownovo gibanje
% Matej Rojec

\begin{document}
\title{Brownovo gibanje}
\author{Matej Rojec}
\date{}
\maketitle

Brownovo gibanje (več v \cite{karatzas1991brownian}) je intuitivno slučajen proces, % Sklic na knjigo
ki predstavlja naključno gibanje delcev v mediju.
    
    % Slika: PerrinPlot2.pdf, velikost scale=0.3
    % Napis pod sliko: 
    % Reprodukcija slike iz Jean Baptiste Perrin, \emph{Mouvement brownien et réalité moléculaire}, Ann. de Chimie et de Physique (VIII) 18, 5-114, 1909
\begin{figure}[ht]
    \centering
    \includegraphics[scale=0.3]{PerrinPlot2.pdf}
    \caption{Reprodukcija slike iz Jean Baptiste Perrin,
    \emph{Mouvement brownien et réalité moléculaire},
    Ann. de Chimie et de Physique (VIII) 18, 5--114, 1909}
\end{figure}{}
    % Začetek definicije
    \begin{definition}
Standardno Brownovo gibanje \(\{B_t\}_{t\geq 0}\) je slučajen proces z naslednjimi lastnostmi:
\begin{enumerate}
    \item \(B_0=0\).
    \item Prirastki \(B_{t_n}-B_{t_{n-1}}, B_{t_{n-1}}-B_{t_{n-2}}, \ldots, B_{t_2}-B_{t_1}, B_{t_1}-B_{t_0}\) so neodvisne slučajne spremenljivke, za vsak \(t_0\leq t_1\leq\cdots\leq t_n\).
    \item Za vsak \(t\geq 0\) in \(h>0\) velja \(B_{t+h}-B_t\sim\mathcal{N}(0,h)\).
    \item Funkcija \(t\mapsto B_t\) je zvezna skoraj gotovo.
\end{enumerate}
\end{definition}
    % Konec definicije
    
    Preden zapišemo izrek, definirajmo še pojem časa ustavljanja.
   
    % Začetek definicije
    \begin{definition}
Slučajna spremenljivka \(\tau\) na verjetnostnem prostoru
\((\Omega,\mathcal{F},P)\) z vrednostmi v
\(\mathbb{R}^{+}\) je čas ustavljanja glede na filtracijo
\((\mathcal{F}_t)_{t\in T}\), če velja:
\[
\forall t\in T:\quad \{\tau\leq t\}\in\mathcal{F}_t.
\]
\end{definition}
    % Konec definicije
    
    Zdaj lahko zapišemo izrek~\ref{thm:stopped_brownian}.
    
    % Začetek izreka
    \begin{theorem}\label{thm:stopped_brownian}
Naj bo \(\{B_t\}_{t\geq 0}\) (standardno) Brownovo gibanje,
\(\tau\) čas ustavljanja glede na
\((\mathcal{F}_t)_{t\geq 0}\) in naj velja
\(P[\tau<\infty]=1\).
Potem je tudi proces
\[
\hat{B}:=\{B_{\tau+t}-B_\tau\mid t\geq 0\}
\]
(standardno) Brownovo gibanje in neodvisen od
\(\mathcal{F}_\tau\).
\end{theorem}
    % Konec izreka
  \bibliographystyle{plain}
\bibliography{knjiga}  
\end{document}